\documentclass{article}
\usepackage[utf8]{utf8}
\usepackage[spanish]{babel}
\usepackage{amsmath, amssymb, amsthm}
\usepackage{geometry}
\geometry{margin=1in}

\title{Prerrequisitos para Operadores de Proyección, Hermiticidad y Bases Ortonormales}
\author{Tutor de Física y Matemáticas}
\date{\today}

\begin{document}

\maketitle

Para comprender adecuadamente los \textbf{operadores de proyección}, \textbf{demostrar la hermiticidad de un operador} y \textbf{construir una base ortonormal}, es necesario dominar los siguientes conceptos del álgebra lineal y la notación de Dirac (braket):

\section{Prerrequisitos para Operadores de Proyección}

Para entender cómo opera un proyector (por ejemplo, $P = \sum_{i} |e_i\rangle\langle e_i|$) sobre un espacio vectorial, se requiere dominar:

\begin{itemize}
    \item \textbf{Propiedad de Idempotencia}: Comprender la condición estructural:
    \[
    P^2 = P
    \]
    \item \textbf{Notación Braket (Representación de Operadores)}:
    \begin{itemize}

        \item Entender el \textbf{producto exterior} $|u\rangle\langle v|$, el cual actúa sobre un ket $|x\rangle$ mediante la contracción $\langle v|x\rangle$, produciendo un nuevo ket proporcional a $|u\rangle$.
        \item La \textbf{relación de clausura o cierre}:
        \[
        \sum_i |e_i\rangle\langle e_i| = I
        \]
        necesaria para expresar el operador identidad o proyectores sobre subespacios.
    \end{itemize}
    \item \textbf{Descomposición en Suma Directa y Ortogonalidad}:
    \begin{itemize}
        \item Descomposición de un espacio $V$ como suma directa $V = M \oplus M^\perp$.
        \item Saber separar un vector en sus componentes pertenecientes a un subespacio $M$ y a su complemento ortogonal $M^\perp$:
        \[
        |v\rangle = P|v\rangle + (I - P)|v\rangle
        \]
    \end{itemize}
\end{itemize}

\section{Prerrequisitos para Demostrar que un Operador es Hermítico}

Para demostrar que un operador $A$ cumple $A^\dagger = A$, se debe conocer:

\begin{itemize}
    \item \textbf{Reglas de la Operación Adjunta ($\dagger$) en Notación Braket}:
    \begin{itemize}
        \item Cambiar kets por bras y viceversa: $(|v\rangle)^\dagger = \langle v|$ y $(\langle w|)^\dagger = |w\rangle$.
        \item Conjugación de escalares: $(c|v\rangle)^\dagger = c^*\langle v|$.
        \item Inversión del orden en los productos: $(AB)^\dagger = B^\dagger A^\dagger$ y $(|u\rangle\langle v|)^\dagger = |v\rangle\langle u|$.
    \end{itemize}
    \item \textbf{Propiedades del Producto Escalar Hermítico}:
    \begin{itemize}
        \item Hermiticidad: $\langle u|v\rangle = \langle v|u\rangle^*$.
        \item Definición del adjunto vía producto escalar:
        \[
        \langle w|A^\dagger|v\rangle = \langle Aw|v\rangle = \langle v|A|w\rangle^*
        \]
    \end{itemize}
    \item \textbf{Identidades de Conmutación y Álgebra de Operadores}:
    \begin{itemize}
        \item Propiedades del conmutador:
        \[
        [A, B] = AB - BA
        \]
    \end{itemize}
\end{itemize}

\section{Prerrequisitos para Crear una Base Ortonormal}

Para construir o verificar una base ortonormal $\{|u_i\rangle\}$, se necesita saber:

\begin{itemize}
    \item \textbf{Independencia Lineal y Concepto de Base}: Comprobar si un conjunto de vectores es linealmente independiente.
    \item \textbf{Condición de Ortonormalidad}:
    \begin{itemize}
        \item Delta de Kronecker:
        \[
        \langle u_i|u_j\rangle = \delta_{ij} = 
        \begin{cases} 
        1 & \text{si } i = j \\ 
        0 & \text{si } i \neq j 
        \end{cases}
        \]
    \end{itemize}
    \item \textbf{Algoritmo de Ortonormalización de Gram-Schmidt}:
    \begin{itemize}
        \item Saber cómo sustraer recursivamente las proyecciones sobre los vectores previa y ortonormalizados:
        \[
        |w_n\rangle = |v_n\rangle - \sum_{j=1}^{n-1} \langle u_j|v_n\rangle |u_j\rangle, \quad |u_n\rangle = \frac{|w_n\rangle}{\||w_n\rangle\|}
        \]
    \end{itemize}
    \item \textbf{Normalización de Vectores}:
    \begin{itemize}
        \item Cálculo de la norma a partir del producto escalar:
        \[
        \|v\| = \sqrt{\langle v|v\rangle}
        \]
    \end{itemize}
\end{itemize}

\end{document}